\doublespacing % Do not change - required

\chapter{Method and Implementation}
\label{ch:method-implementation}

%%%%%%%%%%%%%%%%%%%%%%%%%%%%%%%%%%%%%%%
% IMPORTANT
\begin{spacing}{1}
\minitoc
\end{spacing}
\thesisspacing
%%%%%%%%%%%%%%%%%%%%%%%%%%%%%%%%%%%%%%%

\section{Design Rationale}

The design follows from three limitations in the related work. Mean aggregation
tolerates heterogeneous benign clients, but a large or coordinated
model-poisoning update enters the average with linear influence.
Population-centred robust rules, including centred clipping, reduce this risk
by comparing clients with a common centre, but that centre can be misleading
under non-IID data. Temporal methods show that history can stabilise learning,
but the historical reference should reflect each client's own behaviour rather
than one global pattern shared by all clients.

These observations determine the HRAC components. Like centred clipping, HRAC
uses clipping to bound influence rather than to make Byzantine updates vanish.
The global median norm cap limits large-scale malicious updates before they
enter the aggregate. The history-residual clipping step then allows each client
to affect the round only through a bounded residual relative to its own recent
behaviour, so a client is not penalised only because it is far from the
population centre. The \(\nu_i\) statistic tracks residual movement over time,
and the weight penalty suppresses persistently elevated movement
gradually before the weighted aggregation. HRAC is therefore not a hard
detector. It is an influence-limiting rule designed for settings where benign
heterogeneity and malicious update deviations can coexist.
\section{Problem Setup and Threat Model}
\label{sec:problem-threat-model}

Consider a federated learning system with $N$ clients and a global
server. At communication round $t$, each participating client computes a local
training message $\Delta_i^t \in \mathbb{R}^d$ and sends it to the server. As
in Chapter~\ref{ch2}, $\Delta_i^t$ is treated as a gradient-like descent
direction, so the global model is updated by subtracting the final aggregate.
The server applies an aggregation rule
\[
    \operatorname{Agg}(\Delta_1^t,\ldots,\Delta_N^t)
\]
to form a single update $g_t$, which is then used to update the global model.

The client population is divided into an unknown benign set $\mathcal{G}$ and a
Byzantine set $\mathcal{B}$, with $|\mathcal{B}|=b$ and $|\mathcal{G}|=N-b$.
Clients in $\mathcal{G}$ follow the training protocol, but their updates may be
heterogeneous because their local data are non-IID. Clients in $\mathcal{B}$
may send arbitrary or strategically crafted messages. In the experiments, this
includes both direct model manipulation attacks, such as bit flipping, ALIE,
IPM, MSA, and HisMSA, and protocol-compliant data poisoning through label
flipping.

The objective is not to perfectly classify clients as benign or malicious.
Instead, the goal is to limit the influence of harmful updates while preserving
useful contributions from benign but heterogeneous clients. This distinction is
important under non-IID data: a benign client may be far from the population
centre for legitimate statistical reasons, so hard outlier removal can damage
utility.

HRAC is designed for this setting with two assumptions. First, extreme
round-wise update magnitudes should be bounded before aggregation. Second,
client behaviour should be compared primarily with the client's own historical
trajectory rather than only with a global population centre.

\section{Aggregation Pipeline}

In communication round $t$, the server receives flattened client messages
$\{\Delta_i^t\}_{i=1}^{N}$, with $\Delta_i^t \in \mathbb{R}^d$. HRAC first
limits extreme magnitudes using a median-based cap, then processes each client
relative to its own residual history. The clipped residuals are used to score
cross-round movement through \(\nu_i\). After a warm-up period, the server
turns these scores into soft weights, forms the weighted aggregate, and then
updates the per-client histories.

Algorithm~\ref{alg:hrac} summarises the complete server-side HRAC procedure.
The rest of this section gives the definitions and design rationale behind each
stage.

\begin{algorithm}[p]
\caption{History-Residual Adaptive Clipping (HRAC)}
\label{alg:hrac}
\begin{algorithmic}[1]
\Require messages $\{\Delta_i^t\}_{i=1}^{N}$, client ids, states
$(b_i,\mu_i,\nu_i,p_i,\bar r_i^{\mathrm{prev}},\bar m)$, parameters
$c_g,c,\rho_g,\rho_b,\rho_\mu,\rho_\nu,\rho_p,T_\nu,\alpha,w_{\max}$
\Ensure aggregated update $g_t$
\State $m_t \gets \operatorname{med}_{j}\|\Delta_j^t\|_2$
\State if $\bar m$ is uninitialised, set $\bar m\gets m_t$; \quad
$B_t \gets c_g\max(\bar m,s_{\min})$
\For{each client $i$}
    \State $\hat{\Delta}_i^t \gets
    \min(1,B_t/(\|\Delta_i^t\|_2+\epsilon))\Delta_i^t$
    \If{$i$ is a new client}
        \State $b_i\gets\hat\Delta_i^t$; \quad
        $\mu_i,\nu_i\gets \min(\max(m_t,\mu_{\min}),2B_t)$
        \State $p_i\gets0$; \quad $\bar r_i^{\mathrm{prev}}\gets0$
        \State $\tilde{\Delta}_i^t\gets \hat{\Delta}_i^t$
    \Else
        \State $r_i^t\gets \hat{\Delta}_i^t-b_i$; \quad
        $\tau_i^t\gets c\mu_i+\epsilon$
        \State $\bar r_i^t\gets
        \min(1,\tau_i^t/(\|r_i^t\|_2+\epsilon))r_i^t$
        \State $\tilde{\Delta}_i^t\gets b_i+\bar r_i^t$
    \EndIf
\EndFor
\If{$t\leq T_\nu$}
    \State $w_i^t\gets 1/N$ for all clients
\Else
    \State $\bar\nu^t\gets N^{-1}\sum_{j=1}^{N}\nu_j$
    \For{each client $i$}
        \State $e_i^t\gets \mathbb{I}[\nu_i>0.9]\max(\nu_i-\bar\nu^t,0)$
        \State $p_i\gets \rho_p p_i+(1-\rho_p)e_i^t$
        \State $\hat w_i^t\gets 1/(1+\alpha p_i)$
    \EndFor
    \State normalise $\{\hat w_i^t\}$ and cap weights by $w_{\max}$
\EndIf
\State $g_t\gets \sum_{i=1}^{N} w_i^t\tilde{\Delta}_i^t$
\For{each non-new client $i$}
    \State $\tilde{\Delta}_{i,\mathrm{safe}}^t\gets\tilde{\Delta}_i^t$
    \State $b_i\gets \rho_b b_i+(1-\rho_b)\tilde{\Delta}_{i,\mathrm{safe}}^t$
    \State $\mu_i\gets \rho_\mu\mu_i+(1-\rho_\mu)\|\bar r_i^t\|_2$
    \State $d_i^t\gets \|\bar r_i^t-\bar r_i^{\mathrm{prev}}\|_2$
    \State $\nu_i\gets \rho_\nu\nu_i+(1-\rho_\nu)d_i^t$; \quad
    $\bar r_i^{\mathrm{prev}}\gets \bar r_i^t$
\EndFor
\State $\bar m\gets \rho_g\bar m+(1-\rho_g)m_t$
\State \Return $g_t$
\end{algorithmic}
\end{algorithm}

\section{Core Components}

\subsection{Global Median Norm Cap}

The first defence layer prevents a single large update from dominating the
aggregate. HRAC smooths the round-wise median update norm with an EMA before
forming the clipping threshold. Let
\[
    m_t = \operatorname{med}_{1\leq j\leq N}\|\Delta_j^t\|_2,
    \qquad
    \bar m_t = \rho_g \bar m_{t-1} + (1-\rho_g)m_t .
\]
The clipping threshold used in round $t$ is formed from the stored scale at the
start of the round,
\[
    B_t = c_g\max\{\bar m_{t-1},s_{\min}\},
\]
where \(s_{\min}=10^{-3}\) in the implementation. The first round warm-starts
\(\bar m\) from the current median norm. Thus,
the current round contributes to the next round's cap rather than immediately
changing its own threshold. This makes the cap predictable with respect to the
current batch of submitted messages and reduces sensitivity to one-round norm
spikes or drops. Each message is clipped by
\begin{equation}
    \hat{\Delta}_i^t
    =
    \min\!\left(1,\frac{B_t}{\|\Delta_i^t\|_2+\epsilon}\right)\Delta_i^t .
\end{equation}
The implementation uses $c_g=3.0$ and $\rho_g=0.95$ by default. The cap remains
median-based, so it does not require knowing which clients are malicious and is
less sensitive to extreme Byzantine norms than a mean-based scale estimate; the
additional EMA smoothing makes the cap less volatile across communication
rounds.

\subsection{Per-Client History-Residual Clipping}

For each client $i$, HRAC stores a historical centre $b_i^t$. This centre is an
EMA of safe past messages and represents the client's own typical update path,
not a population-wide prototype. The residual is
\begin{equation}
    r_i^t = \hat{\Delta}_i^t - b_i^t .
\end{equation}
The residual threshold is adaptive:
\begin{equation}
    \tau_i^t = c\mu_i^t + \epsilon ,
\end{equation}
where $\mu_i^t$ is an EMA scale statistic and $c=2.5$ by default. The clipped
residual is
\begin{equation}
    \bar r_i^t
    =
    \min\!\left(1,\frac{\tau_i^t}{\|r_i^t\|_2+\epsilon}\right)r_i^t .
\end{equation}
The processed message is then reconstructed as
\begin{equation}
    \tilde{\Delta}_i^t = b_i^t + \bar r_i^t .
\end{equation}
This reconstructed message is used directly for aggregation after residual
clipping.
In the reported runs, the same reconstructed message is also used to update the
client history centre, so
\[
    \tilde{\Delta}_{i,\mathrm{safe}}^t=\tilde{\Delta}_i^t .
\]

\subsection{Separated Aggregation and Statistics Paths}

The implementation distinguishes the aggregation path from the statistics path.
Aggregation uses the globally clipped update $\hat{\Delta}_i^t$ and its
processed version $\tilde{\Delta}_i^t$. The statistics path can lift extremely
small-norm updates to the median norm when computing $\mu_i$ and $\nu_i$. This
prevents a near-zero update from artificially collapsing a client's history
statistics. The aggregation value itself is not lifted in this case, so the
actual descent direction remains faithful to the received message after global
clipping.

\subsection{EMA History Updates}

After aggregation, HRAC updates the client histories. The aggregation-path
centre is
\begin{equation}
    b_i^{t+1}
    =
    \rho_b b_i^t + (1-\rho_b)\tilde{\Delta}_{i,\mathrm{safe}}^t ,
\end{equation}
where $\tilde{\Delta}_{i,\mathrm{safe}}^t$ is the message used in the history
update.
For the reported runs, this is the same reconstructed message used in the
aggregate. The default value is $\rho_b=0.98$.

The residual scale statistic is updated as
\begin{equation}
    \mu_i^{t+1}
    =
    \rho_\mu \mu_i^t + (1-\rho_\mu)\|\bar r_i^t\|_2 ,
\end{equation}
with a small floor $\mu_{\min}$ for numerical stability. The default is
$\rho_\mu=0.95$.

The change statistic $\nu_i$ is updated from the movement of the clipped
residual:
\begin{equation}
    d_i^t = \|\bar r_i^t - \bar r_i^{t-1}\|_2 ,
    \qquad
    \nu_i^{t+1}
    =
    \rho_\nu \nu_i^t + (1-\rho_\nu)d_i^t ,
\end{equation}
again with a small floor $\nu_{\min}$. The class default is $\rho_\nu=0.95$.
In the reported experiment driver, this value is set to $\rho_\nu=0.87$ so that
the residual-change statistic reacts slightly faster to persistent movement.

\subsection{Soft Weighting from the \texorpdfstring{$\nu$}{nu} Statistic}

The weighting rule is deliberately delayed. HRAC does not immediately
down-weight clients from the first round. For
iterations $t \leq T_{\nu}$, it uses uniform weights. In the implementation,
$T_{\nu}=50$ by default. This warm-up allows the per-client statistics to become
meaningful before they affect aggregation.

Once the penalty is active, the server forms a population baseline
\begin{equation}
    \bar{\nu}^t = \frac{1}{N}\sum_{j=1}^{N}\nu_j^t .
\end{equation}
Only clients with sufficiently high $\nu_i^t$ are penalised. In code, the
activation threshold is $0.9$, and the excess score is
\begin{equation}
    e_i^t =
    \begin{cases}
        \max(\nu_i^t-\bar{\nu}^t,0), & \nu_i^t>0.9,\\
        0, & \nu_i^t\leq 0.9 .
    \end{cases}
\end{equation}
The excess score is itself smoothed before being used for aggregation:
\begin{equation}
    p_i^t = \rho_p p_i^{t-1} + (1-\rho_p)e_i^t .
\end{equation}
The implementation uses \(\rho_p=0.90\). This extra EMA prevents a single
large residual-change measurement from causing an abrupt weight change.
The raw weight is then
\begin{equation}
    \hat w_i^t = \frac{1}{1+\alpha p_i^t},
\end{equation}
where $\alpha=5.0$ by default. Raw weights are normalised to sum to one. A
maximum per-client weight cap, $w_{\max}=0.30$, is then applied iteratively with
renormalisation to avoid a small number of clients dominating the round.
More explicitly, let \(q_i=\hat w_i^t/\sum_k\hat w_k^t\). If all
\(q_i\le w_{\max}\), then \(w_i^t=q_i\). Otherwise, all clients with
\(q_i>w_{\max}\) are fixed at \(w_{\max}\); the remaining probability mass is
renormalised over the uncapped clients in proportion to their current \(q_i\)
values. This fix-and-renormalise step is repeated until no weight exceeds
\(w_{\max}\). The resulting weights are the final aggregation weights used in
\(g_t\). In the convergence analysis, quantities such as the Byzantine weight
mass \(\beta_t\) are defined after this capping step.

The final aggregate is
\begin{equation}
    g_t = \sum_{i=1}^{N} w_i^t\tilde{\Delta}_i^t .
\end{equation}

\section{Software Integration and Logging}
\label{sec:software-integration}

HRAC is implemented as the \texttt{C\_HRAC} aggregation class in the same
centralised aggregation interface used by the baseline rules. The training loop
therefore remains the same across mean aggregation, FABA, centred clipping,
LFighter, and HRAC. The same experiment driver also constructs the attack
object, including label flipping, bit flipping, ALIE, IPM, MSA, and HisMSA, so
the method is evaluated under matched training settings rather than through
attack-specific code.

The ablation experiments use controlled variants of the same HRAC pipeline by
removing the global median cap, per-client residual clipping, $\nu$-based
weighting, or both history-based stages. HRAC also records diagnostic quantities
at a fixed logging interval, including residual thresholds, $\mu$ and $\nu$
statistics, pre- and post-clipping norms, normalised weights, and aggregate
norms. Accuracy, loss, metadata, and diagnostic logs are stored through the
common cache interface and reused by the plotting scripts in
Chapter~\ref{ch:experiments-results}.

\section{Hyperparameters}

The main HRAC hyperparameters are:
\begin{itemize}[leftmargin=*]
    \item $\rho_b=0.98$: EMA decay for the client history centre.
    \item $\rho_\mu=0.95$: EMA decay for the residual scale statistic.
    \item $\rho_\nu=0.87$ in the experiment driver: EMA decay for the
    residual-change statistic. The class default is $0.95$, but the reported
    runs use a slightly faster $\nu$ update.
    \item $c=2.5$: multiplier for the per-client residual clipping threshold.
    \item $c_g=3.0$: multiplier for the EMA-smoothed global median norm cap.
    \item $\rho_g=0.95$: EMA decay for the global median-norm scale.
    \item $T_\nu=50$: first iteration after which $\nu$-based weighting is used.
    \item $\alpha=5.0$: strength of the $\nu$ excess penalty.
    \item $\rho_p=0.90$: EMA decay for the smoothed $\nu$-penalty score.
    \item $w_{\max}=0.30$: maximum normalised weight allowed for one client.
\end{itemize}

The algorithm is intentionally not a hard detector. Its most sensitive
parameters are $c$, which controls how much residual deviation is allowed before
clipping, and $\alpha$, which controls how sharply elevated $\nu$ values reduce
weights.

\section{Convergence Statement}
\label{sec:convergence-statement}

The empirical evaluation in Chapter~\ref{ch:experiments-results} is the primary
evidence for HRAC's practical robustness. The method can also be analysed
through a modular convergence argument. The analysis treats HRAC as an
aggregation rule that produces \(g_t\), compares this aggregate with the honest
conditional centre, and inserts the resulting aggregation error into the
standard descent argument for smooth non-convex optimisation. The lagged
EMA-smoothed median cap makes the clipping radius predictable before the current
round's stochastic updates are observed. Under additional admissibility
conditions on the honest conditional centre and residual scale, this allows the
clipped-moment lemma of Sadiev et
al.~\cite[Lemma~5.1]{sadiev2023highprob} to be applied.

Under bounded honest updates, an honest-majority condition for the median-based
scale, bounded honest clipping distortion, bounded honest-weight distortion, and
bounded residual Byzantine weight mass, this conditional analysis shows that
HRAC converges to a first-order stationary neighbourhood in the averaged-gradient
sense. In
particular, for a sufficiently small constant step size \(\eta\), the averaged
squared gradient norm satisfies a bound of the form
\begin{equation}
    \frac{1}{T}\sum_{t=0}^{T-1}\mathbb{E}\|\nabla f(x_t)\|^2
    \le
    \frac{2(f(x_0)-f^\star)}{\eta T}
    +
    R_{\mathrm{HRAC}},
\end{equation}
where \(R_{\mathrm{HRAC}}\) collects the honest stochastic variance,
non-IID/bias term, global clipping distortion, residual clipping distortion,
and the remaining Byzantine aggregation error. When the logged dynamics exhibit
a late-stage group-level \(\nu\)-gap, the weighting rule can reduce the
residual Byzantine mass appearing in this error term.

This result should be read as a conditional bounded-influence stability
statement rather than as a claim of exact convergence to a zero-gradient point
under arbitrary attacks. The proof does not establish universal Byzantine
identification. Persistent heterogeneity and nonzero Byzantine mass create a
nonzero error floor. The full proof, including the aggregation-error
decomposition and the explicit neighbourhood constant, is given in
Appendix~\ref{app:convergence-analysis}.
